\documentclass[answers,12pt,addpoints]{exam} 
\usepackage{import}

\import{C:/Users/prana/OneDrive/Desktop/MathNotes}{style.tex}

% Header
\newcommand{\name}{Pranav Tikkawar}
\newcommand{\course}{Spatial Statistics}
\newcommand{\assignment}{Homework 5.3}
\author{\name}
\title{\course \ - \assignment}

\begin{document}
\maketitle

\begin{questions}
\question 6.5

\begin{solution}
    \textbf{(a)}
    I believe that I am not super amazing at viewing log-log plots, but the slope of the line seems to be $\frac{1}{3}$,

    When plotting it in desmos, I got the following graph:
    \begin{center}
        \includegraphics[scale=0.3]{img/desmos.png}
    \end{center}

    I also tried combinations of log slopes with the points and the values of the slopes fluctuated a lot between .3 and .5, but I think the best fit is $\frac{1}{3}$.

    \textbf{(b)}
    The paramertization of this model should be given by the first odd power term in the Taylor expansion of the function. The taylor expansion of $\theta_1 e^{-r/ \theta_2} (1+r/\theta_2)$ is given by:
    \begin{align*}
        \theta_1 e^{-r/ \theta_2} (1+r/\theta_2) &= \theta_1 \left( 1 - \frac{r}{\theta_2} + \frac{r^2}{2\theta_2^2} - \frac{r^3}{6\theta_2^3} + \frac{r^4}{24\theta_2^4} \cdots \right) \left( 1 + \frac{r}{\theta_2} \right) \\
        &= \theta_1 \left( 1 - \frac{r^2}{2\theta_2^2} + \frac{r^3}{3\theta_2^3} + \cdots \right) \\
        &= \theta_1 + \frac{\theta_1 r^2}{2\theta_2^2} - \frac{\theta_1 r^3}{6\theta_2^3} + \cdots
    \end{align*}
    Thus the first odd power term is $-\frac{\theta_1 r^3}{3\theta_2^3}$, so the parameterization of this model should be $\phi_1 = \frac{\theta_1}{\theta_2^3}$ and the other $\phi_2 = \frac{1}{\theta_2}$. 

\textbf{(c)}
Considering the $\theta_1 e^{-r/\theta_2} (1+r/\theta_2)$ model, we
reparameterize it in terms of $\phi_1$ and $\phi_2$ as follows:
\[
\theta_1 e^{-r/\theta_2} (1+r/\theta_2)
= \frac{\phi_1}{\phi_2^{3}} e^{-r \phi_2} (1+r \phi_2).
\]

We can also observe the Taylor expansion of this function as follows:
\begin{align*}
\frac{\phi_1}{\phi_2^{3}} e^{-r \phi_2} (1+r \phi_2)
&= \frac{\phi_1}{\phi_2^{3}}
   \left( 1 - r \phi_2 + \frac{r^2 \phi_2^2}{2}
          - \frac{r^3 \phi_2^3}{6} + \cdots \right)
   \left( 1 + r \phi_2 \right) \\
&= \frac{\phi_1}{\phi_2^{3}}
   \left( 1 - \frac{r^2 \phi_2^2}{2}
          + \frac{r^3 \phi_2^3}{3} + \cdots \right) \\
&= \frac{\phi_1}{\phi_2^{3}} - \frac{\phi_1 r^2}{2\phi_2}
   + \frac{\phi_1 r^3}{3} + \cdots .
\end{align*}

For our data $Y \sim N(\beta \mathbf{1}, \Sigma(\phi_1, \phi_2))$,
where $\beta$ is an unknown mean, we can use contrasts to get the
REML estimates of $\phi_2$. Define $Z = C^\top Y$ where $C$ is a
contrast matrix such that $C^\top \mathbf{1} = 0$ and
$C^\top C = I$. Then we have that
$Z \sim N(0, C^\top \Sigma(\phi_1, \phi_2) C)$.

For the REML log-likelihood (up to an additive constant) we can use
the following formula:
\[
\ell_{R}(\phi_1, \phi_2)
= -\frac{1}{2} \log \bigl|C^\top \Sigma(\phi_1, \phi_2) C\bigr|
  -\frac{1}{2} Z^\top
     \bigl(C^\top \Sigma(\phi_1, \phi_2) C\bigr)^{-1} Z.
\]

Next, we observe $C^\top \Sigma(\phi_1, \phi_2) C$ as
$\phi_2 \to 0$. Let us first define our covariance matrix
$\Sigma(\phi_1, \phi_2)$ as follows:
\[
\Sigma(\phi_1, \phi_2) = \bigl( K(r_{ij}) \bigr)_{i,j}.
\]
For convenience we define
$(R_{\phi_2})_{ij} = e^{-r_{ij} \phi_2} (1+r_{ij} \phi_2)$, then we
can write $\Sigma(\phi_1, \phi_2) = \dfrac{\phi_1}{\phi_2^3} R_{\phi_2}$.
Therefore the contrast covariance matrix can be written as
\[
C^\top \Sigma(\phi_1, \phi_2) C
= \frac{\phi_1}{\phi_2^3} C^\top R_{\phi_2} C.
\]

Now observing $R_{\phi_2}$ as a matrix Taylor expansion, we have that
\[
R_{\phi_2}
= \mathbf{1}\mathbf{1}^\top
  - \frac{\phi_2^2}{2} R^{(2)}
  + \frac{\phi_2^3}{3} R^{(3)} + \cdots,
\]
where $R^{(2)}$ and $R^{(3)}$ are the matrices of the second and third
powers of the distances between the points. Thus we have that
\[
C^\top R_{\phi_2} C
= C^\top \mathbf{1}\mathbf{1}^\top C
  - \frac{\phi_2^2}{2} C^\top R^{(2)} C
  + \frac{\phi_2^3}{3} C^\top R^{(3)} C + \cdots.
\]
Since $C^\top \mathbf{1} = 0$, we lose the first term and we have that
\[
C^\top R_{\phi_2} C
= - \frac{\phi_2^2}{2} C^\top R^{(2)} C
  + \frac{\phi_2^3}{3} C^\top R^{(3)} C + \cdots.
\]
Thus we have that
\[
C^\top \Sigma(\phi_1, \phi_2) C
= - \frac{\phi_1}{2\phi_2} C^\top R^{(2)} C
  + \frac{\phi_1}{3} C^\top R^{(3)} C + \cdots.
\]

Next we observe the leading term to understand the behavior of the
REML estimate as $\phi_2 \to 0$. Let the locations of the points be
given by the rows of a matrix $X \in \mathbb{R}^{m \times k}$.
If we consider that $r_{ij} = \| x_i - x_j\|$ then
$R^{(2)}_{ij} = r_{ij}^2$. By the identity
$r_{ij}^2 = \|x_i\|^2 + \|x_j\|^2 - 2 x_i^\top x_j$, we can write
$R^{(2)}$ as follows, with $s$ the $m$–vector of squared norms of the
rows of $X$:
\[
R^{(2)} = \mathbf{1} s^\top + s \mathbf{1}^\top - 2 X X^\top.
\]
When we conjugate this with $C$, we have that
\[
C^\top R^{(2)} C
= - 2 C^\top X X^\top C.
\]
Let us define $B = X^\top C$, then we have that
\[
C^\top R^{(2)} C = - 2 B^\top B.
\]
We know that $B^\top B$ is positive semidefinite and has rank at most
$k$. Thus $C^\top R^{(2)} C$ is negative semidefinite and has rank at
most $k$. We can further say that $C^\top R^{(2)} C$ has exactly rank
$k$ since $C$ has full column rank $m-1$ and $X$ has full column rank
$k$. Thus all $k$ nonzero eigenvalues of $C^\top R^{(2)} C$ are
negative.

If $\lambda_1, \ldots, \lambda_k$ are the nonzero eigenvalues of
$C^\top R^{(2)} C$, then for small $\phi_2$
\[
C^\top \Sigma(\phi_1, \phi_2) C
\approx - \frac{\phi_1}{2\phi_2} C^\top R^{(2)} C,
\]
so the $k$ nonzero eigenvalues of $C^\top \Sigma(\phi_1, \phi_2) C$
are approximately given by $\sim c_i \dfrac{\phi_1}{\phi_2}$ where
$c_i = -\lambda_i/2 > 0$. Hence
\[
\log \bigl| C^\top \Sigma(\phi_1, \phi_2) C \bigr|
= \sum_{i=1}^{m-1} \log \lambda_i(\phi_2)
\sim k \log\!\left(\frac{\phi_1}{\phi_2}\right) + O(1),
\]
and therefore
\[
-\frac{1}{2} \log \bigl| C^\top \Sigma(\phi_1, \phi_2) C \bigr|
\sim -\frac{k}{2} \log\!\left(\frac{\phi_1}{\phi_2}\right) + O(1).
\]

For the quadratic term in $\ell_R$, diagonalize
$C^\top \Sigma(\phi_1,\phi_2) C$ and write $Z$ in the corresponding
eigenbasis. Then
\[
-\frac{1}{2} Z^\top
   \bigl(C^\top \Sigma(\phi_1,\phi_2) C\bigr)^{-1} Z
= -\frac{1}{2} \sum_{i=1}^{m-1}
   \frac{\tilde z_i^2}{\lambda_i(\phi_2)}.
\]
For the $k$ large eigenvalues we have
$\lambda_i(\phi_2) \sim c_i \phi_1 \phi_2^{-1}$, so
$1/\lambda_i(\phi_2) = O(\phi_2) \to 0$ as $\phi_2 \to 0$, and these
contributions vanish. The remaining $m-k-1$ eigenvalues stay bounded
away from zero and infinity, so the corresponding terms converge to a
finite limit. Thus the quadratic term remains finite as
$\phi_2 \to 0$.

Combining these results, as $\phi_2 \to 0$ the determinant term in
$\ell_R(\phi_1,\phi_2)$ behaves like
$-\frac{k}{2} \log(\phi_1/\phi_2) \to -\infty$, while the quadratic
term stays finite. Hence, for this parametrization and an unknown
constant mean, the restricted log-likelihood does \emph{not} tend to a
finite value as the inverse range parameter $\phi_2$ tends to $0$.
\end{solution}


\question 6.13

\begin{solution}
\[
K(r) = \frac{\phi_1}{\phi_2}\exp(-\phi_2 r),
\]

To study how often the REML estimate of $\phi_2$ hits the boundary at
$0$, I ran the following simulation. I fixed $\phi_1 = 1$ and
considered sample sizes
\[
m \in \{5, 10, 50, 100, 200\}
\]
and true inverse ranges
\[
\phi_2 \in \{0.2, 0.5, 1, 2, 5\}.
\]
For each pair $(m,\phi_2)$ I repeated the steps below $N = 200$ times:
\begin{enumerate}
\item Generate $m$ locations independently from $\mathrm{Unif}([0,1]^2)$
      in the unit square.
\item Form the $m\times m$ covariance matrix $\Sigma(\phi_1,\phi_2)$
      with entries
      $\Sigma_{ij} = K(r_{ij})
      = (\phi_1/\phi_2)\exp(-\phi_2 r_{ij})$.
\item Simulate one realization
      $Y \sim N(\beta\mathbf{1},\Sigma(\phi_1,\phi_2))$ with $\beta=0$.
\item Fit the same model by restricted maximum likelihood, using
      \texttt{optimize} to maximize the restricted log-likelihood over
      $\phi_2 \ge 0$ while profiling out $\phi_1$ and the mean.
\item Record whether the REML estimate satisfies
      $\hat\phi_2^{(R)} = 0$. In the code I treat
      $\hat\phi_2^{(R)} < 10^{-4}$ as $0$.
\end{enumerate}

For each $(m,\phi_2)$ I then estimated
$
\hat p(m,\phi_2)
= \frac{1}{N} \sum_{r=1}^N
  \mathbf{1}\{\hat\phi_{2,r}^{(R)} = 0\}.
$
The plot below shows $\hat p(m,\phi_2)$ as a function of $m$ (on a log
scale) for each fixed value of $\phi_2$.

Numerically, for $\phi_2 = 0.2$ I obtained
\[
\hat p(5,0.2) = 0.35,\quad
\hat p(10,0.2) = 0.355,\quad
\hat p(50,0.2) = 0.295,\quad
\hat p(100,0.2) = 0.32,\quad
\hat p(200,0.2) = 0.34,
\]
so the probability that the REML estimate hits the boundary stays
around $0.3$ and does not seem to decrease as $m$ increases. For
$\phi_2 = 0.5$ the estimated probabilities lie between about $0.24$
and $0.45$ over this range of $m$, again clearly bounded away from $0$.
For $\phi_2 = 1$ the probability decreases from about $0.40$ at $m=5$
to about $0.12$ at $m=200$. For larger inverse ranges,
$\phi_2 = 2$ and $\phi_2 = 5$, the probability becomes quite small:
for $\phi_2 = 2$ it drops to roughly $0.02$ by $m=200$, and for
$\phi_2 = 5$ it is essentially $0$ once $m \ge 50$.

Overall, these simulations suggest that for small inverse ranges
(i.e., very smooth, strongly correlated fields) the probability that
$\hat\phi_2^{(R)} = 0$ does not tend to $0$ as $m$ increases. Instead,
for values like $\phi_2 = 0.2$ and $\phi_2 = 0.5$,
$\hat p(m,\phi_2)$ appears to stabilize around a positive level. For
larger inverse ranges, such as $\phi_2 = 2$ or $\phi_2 = 5$, the
probability does decrease toward $0$ as $m$ grows. This is consistent
with the discussion in Section 6.3 and the restricted log-likelihood
plots in Figure 6.11: because the restricted log-likelihood has a
finite limit as $\phi_2 \downarrow 0$ for this parametrization and an
unknown constant mean, there can be a non-negligible chance that the
REML maximizer over $\phi_2 \ge 0$ occurs at the boundary, especially
when the true process has a very large spatial range.

Note that I got help for the code and simulation approach from George and Luke.
\end{solution}

\begin{center}
    \includegraphics[scale=.4]{img/5_3.png}
\end{center}




\end{questions}

\end{document}